% ============================================================
%  SEÇÃO 1 — INTRODUÇÃO
%
%  A introdução deve responder a 4 perguntas:
%    1. Qual é o contexto do problema? (por que isso importa?)
%    2. Qual é o problema de pesquisa? (o que você vai investigar?)
%    3. Quais são os seus objetivos?
%    4. Como este trabalho está organizado?
%
%  Tamanho sugerido: 2 a 4 páginas
% ============================================================

\chapter{INTRODUÇÃO}

% ----------------------------------------------------------
% 1.1 CONTEXTUALIZAÇÃO
% Apresente o cenário em que o problema está inserido.
% Use dados, estatísticas ou referências para situar o leitor.
% ----------------------------------------------------------

\section{Contextualização}

% %->  ESCREVA AQUI a contextualização do seu tema.
%     Responda: Por que este assunto é relevante hoje?
%               Qual é o cenário que motivou sua pesquisa?
%
% Exemplo de início:
% A crescente digitalização dos serviços públicos no Brasil tem
% evidenciado desafios relacionados à eficiência no atendimento
% presencial \cite{autor2022}. Segundo \textcite{autor2023}, ...

[Escreva a contextualização aqui. Cite ao menos 2 referências neste parágrafo.]


% ----------------------------------------------------------
% 1.2 PROBLEMA DE PESQUISA
% Uma pergunta clara, específica e investigável.
% ----------------------------------------------------------

\section{Problema de Pesquisa}

% %->  ESCREVA AQUI seu problema de pesquisa.
%     Deve ser formulado como uma pergunta.
%     Dica: copie do seu formulário de proposta de tema.
%
% Estrutura recomendada:
% Diante do contexto apresentado, este trabalho busca responder
% à seguinte questão: [SUA PERGUNTA DE PESQUISA]?

Diante do contexto apresentado, este trabalho busca responder à seguinte questão de pesquisa: \textbf{[Escreva aqui sua pergunta de pesquisa?]}


% ----------------------------------------------------------
% 1.3 OBJETIVOS
% ----------------------------------------------------------

\section{Objetivos}

\subsection{Objetivo Geral}

% %->  ESCREVA AQUI seu objetivo geral.
%     Deve começar com verbo no infinitivo:
%     Desenvolver / Analisar / Avaliar / Propor / Investigar
%
% Exemplo:
% Desenvolver um sistema web para gerenciamento de filas em
% postos de saúde do município de Tianguá-CE.

[Escreva aqui o objetivo geral. Começa com verbo no infinitivo.]


\subsection{Objetivos Específicos}

% %->  LISTE AQUI seus objetivos específicos (3 a 5 itens).
%     Cada um deve ser uma etapa concreta para atingir o geral.
%     Também começam com verbo no infinitivo.
%
% Exemplo:
% \begin{enumerate}
%   \item Levantar os requisitos funcionais e não funcionais
%         junto aos usuários do sistema;
%   \item Modelar a arquitetura do sistema usando diagramas UML;
%   \item Implementar o sistema utilizando tecnologias web modernas;
%   \item Avaliar a usabilidade por meio de testes com usuários reais.
% \end{enumerate}

\begin{enumerate}
    \item {[Primeiro objetivo específico — etapa 1 para atingir o geral]};
    \item {[Segundo objetivo específico — etapa 2]};
    \item {[Terceiro objetivo específico — etapa 3]};
    \item {[Quarto objetivo específico — etapa 4, se necessário]}.  
\end{enumerate}


% ----------------------------------------------------------
% 1.4 JUSTIFICATIVA
% Por que este trabalho é importante?
% ----------------------------------------------------------

\section{Justificativa}

% %->  ESCREVA AQUI a justificativa do seu trabalho.
%     Responda 3 perguntas:
%     (a) Qual problema concreto este trabalho resolve?
%     (b) Quem se beneficia com este trabalho?
%     (c) Por que agora? Por que em SI?
%
%     Use dados e referências para sustentar seus argumentos.
%     Tamanho sugerido: 1 a 2 parágrafos.

[Escreva aqui a justificativa. Use referências para embasar a importância do seu tema.]


% ----------------------------------------------------------
% 1.5 ESTRUTURA DO TRABALHO
% Explica como o documento está organizado.
% ✅ Preencha após finalizar todas as seções.
% ----------------------------------------------------------

\section{Estrutura do Trabalho}

% %->  Adapte o texto abaixo conforme suas seções reais.

Este trabalho está organizado da seguinte forma: o Capítulo~\ref{cap:referencial} apresenta o referencial teórico, abordando os principais conceitos relacionados ao tema. O Capítulo~\ref{cap:metodologia} descreve a metodologia de pesquisa adotada. O Capítulo~\ref{cap:cronograma} apresenta o cronograma de execução do trabalho.
